% 第八章 展望与结论
\FloatBarrier

\chapter{展望与结论}
\label{ch:outlook}

\FloatBarrier

\section{研究总结}

本研究围绕有机光电存算一体（Compute-in-Memory, CIM）系统的算法--器件协同设计，建立了面向混合模拟/数字推理的行为级仿真与评估流程，并将其用于视觉 Transformer 与大语言模型 KV 缓存两类负载。贯穿全文的核心问题是：当器件具备多级电导、光敏性和低压编程等优势时，这些优势能否在系统级任务精度、跨实例鲁棒性和能效边界上真正保留下来。本文的答案是有条件的：有机光电 CIM 可以在受控的模拟层、足够的 ADC 位宽和分布感知训练协议下保持可用精度，但固定硬件实例训练、低位宽读出、严重漂移和不匹配噪声律会迅速压缩部署空间。

第~\ref{chap:methodology}~章建立了从器件物理参数到系统精度的一阶行为级仿真框架。框架覆盖差分对电导映射、多电平量化、D2D/C2C 噪声、双指数保持衰减、光电前端补偿与硬件感知训练等模块。以 Tiny-ViT-5M 为例，$87.7\%$ 参数被映射到模拟域，覆盖注意力投影与前馈网络中的密集算子；动态注意力、Softmax 和归一化保留在数字域，从而明确了模拟加速的作用范围和上限。

第~\ref{ch:failure_modes}~章和第~\ref{chap:deployment-envelope}~章显示，Transformer 对模拟非理想性的敏感度高于传统 CNN，尺度恢复对维持激活动态范围不可或缺。更关键的是，固定 D2D 掩膜会导致硬件实例过拟合：模型在训练实例上表现良好，却在未见硬件实例上退化至 $10.00\%$。Ensemble HAT 通过在训练轮次边界重采样 D2D 失配图，把器件失配纳入训练分布，使 fresh-instance 精度恢复至 $86.16\pm0.19\%$。该现象构成本文最重要的算法结论：模拟噪声鲁棒性不等同于硬件实例迁移鲁棒性，训练目标必须显式覆盖部署实例分布。

多维物理应力测试进一步给出了部署包络。ADC 分辨率存在 6-bit 悬崖，低于该阈值时训练协议难以挽救近随机退化；越过该悬崖后，D2D 失配成为剩余方差的主要来源。CIFAR-100 等细粒度任务表明，CIFAR-10 上有效的保护排序不能直接外推：CIFAR-100 fresh 全新失配可降至约 $1\%$ 的百类随机基线，top-$30$ D2D 保护也仅恢复至 $23.81\pm1.66\%$。在 CIFAR-10 上，层敏感度与选择性保护实验揭示了长尾风险结构，top-$30$ 敏感层保护在 $10$ 个未见 D2D 实例、每实例 $3$ 次 C2C 蒙特卡洛的复核中达到 $86.04\pm0.95\%$，为混合精度阵列分配、冗余保护和局部校准提供了层级筛选依据。

第~\ref{ch:work2_scope}~章将评估范围扩展至 LLM KV 缓存。Remote~107 早期审计表明，全层模拟 KV 缓存在 8-bit 零噪声下已超过 $10\%$ 困惑度退化门限（17.48 vs. 15.68 PPL），不适合作为主路线；后续 selective-KV claim-lock packet 在 Pythia-410M、WikiText-2 raw-v1 test、context length 512、non-overlapping stride 512、batch size 1 下完成 41/41 个逐行 JSON artifact 归档，数字基线为 PPL=23.30。该锁定包显示 last1/last2 在 D2D=0.02--0.05 下保持在约 18.6--19.0 PPL，last4 升至约 20.4--21.5 PPL，而全层 all24 在 D2D=0.05 时恶化到 $59.055\pm1.610$ PPL。这一结果提示有机 CIM 在 LLM 场景中的合理入口不是全层模拟化，而是选择性缓存、局部模拟化和分层保护。

\FloatBarrier

\section{未来研究方向}

后续视觉任务评估需要从 CIFAR 系列扩展至 ImageNet-1K 等大规模数据集，并纳入 DeiT-S、Swin-T 等更大视觉 Transformer。模型规模提升会同时考验阵列容量、片上带宽、外围读出精度和数字--模拟调度；对于更大参数量的模型，全模型一次性映射并不现实，分层映射、选择性模拟化和多芯片阵列铺排将成为必要路线。

混合精度与选择性保护策略需要从诊断实验走向可实现的电路--算法协同方案。当前 \texttt{freeze\_topK\_d2d} 证明了层敏感度排序能显著缩小全新 D2D 实例下的精度鸿沟，但它仍是推理期诊断。后续应把这一排序转化为具体硬件机制，例如高精度阵列分配、冗余列保护、局部写验证、层级 ADC 位宽配置或数字旁路。只有当保护策略对应到可核算的面积、能耗和编程复杂度时，混合精度结论才能从候选设计变为部署方案。

围绕保持漂移的本地后续复核还提示，保护排序本身也应从静态 D2D 敏感度扩展为时间相关的漂移感知排序。在 Tiny-ViT V4 / CIFAR-100 的两组完整 $10\times3$ retention/protection 复核中，基于漂移向量构造的 drift-aware 排序在 seed789 上相对旧的 D2D 敏感度排序于 $0/1000/10000$~s 分别带来 top-$30$ 的 $+1.38/+1.58/+1.51$~pp 和 top-$42$ 的 $+2.70/+2.83/+2.81$~pp 改进，同时保持 \texttt{fresh\_all\_analog} 基线不变；在更强的 seed456 检查点上，同一排序相对 fresh-all 基线仍保持 top-$30$ 的约 $+2.09/+2.06/+1.94$~pp 和 top-$42$ 的约 $+3.31/+3.31/+3.25$~pp 增益。该结果仍属于本地 provisional 证据，但已表明后续混合精度分配不应只看瞬时 D2D 风险，还应显式考虑保持漂移方向与幅度。更进一步的下一步，是把 drift-aware 排序升级为训练期目标或正则项，而不只是推理期保护启发式。

LLM KV 缓存方向应继续沿选择性终端层和长上下文保持建模推进。TinyLlama-1.1B 与 107-clean Pythia-410M 的阶段性结果共同提示，终端层 KV 缓存比全层缓存更适合有机 CIM 的噪声和保持特性，但不同协议的数字基线必须分开标注，不能把 PPL 数值直接合并。未来需要在具备分组查询注意力的 LLaMA 系列或同类模型上验证该规律，并把上下文长度、KV 生命周期、刷新间隔、光学恢复机制和困惑度退化统一建模。对于 32k 及以上上下文，关键问题不是单次读出精度，而是分钟级会话中保持漂移、刷新能耗和缓存替换策略之间的系统折中。

物理实验闭环是从行为级仿真走向工艺协同设计的关键环节。本文目前主要基于文献先验和参数化模型，尚未与真实有机光电器件的晶圆级测量数据闭环。未来应建立从实验表征到仿真参数的自动化转换管线，将多级编程 I--V 曲线、电导窗口、积分非线性、C2C 方差、跨芯片 D2D 分布、保持衰减轨迹和光响应非线性统一拟合为可直接驱动训练与推理脚本的器件画像。默认双指数保持模型（$\tau_1=140$~ms, $\tau_2=610$~ms, $A_0=0.6$）也应替换为 DNTT、C8-BTBT 及其衍生材料在不同温度与偏置下的实测值；NL\_LTP/NL\_LTD 的梯度缩放近似则应升级为脉冲级写验证动力学模型。

系统能效模型也需要从一阶估算走向外围电路显式建模。现有模型将阵列到数字互连能耗部分归入 SRAM 读写项，尚未单独核算专用布线、数据编排、模拟--数字接口和层间尺度恢复乘法器。更精细的模型应包括分块 ADC/DAC 转换次数、可编程跨阻放大器增益校准、分层自适应 ADC 配置、位置相关 IR 压降、潜行电流和温度漂移的 SPICE 校准。对于长期推理，还需要闭环写验证、周期性刷新、自适应偏置补偿和动态余量缩放。KV 缓存场景中的光学刷新尤其值得单独评估，因为它可能在不消耗电写耐久预算的情况下维持电导态，是有机 OEC-RAM 相对于纯电存储器的重要潜在优势。

当单层阵列无法承载大模型参数或长上下文缓存时，多芯片阵列铺排与片间路由将成为系统瓶颈。未来研究需要处理模拟权重在多个物理芯片间的划分、片间模拟信号完整性、数字控制逻辑的分布式调度与同步，以及光互连在三维堆叠或晶圆级键合中的集成方式。有机光电 CIM 的潜在优势不仅在于存算融合本身，还在于光调控提供的多模态编程和通信接口；这一特点可能支撑从单阵列原型向可扩展系统级平台演进。

\FloatBarrier

\section{结论}

有机光电存算一体技术为突破冯--诺依曼架构能效瓶颈提供了重要路径，但现代深度学习推理到有机阵列的映射并不是简单权重搬迁。动态注意力的数字开销、ADC 量化阈值、保持漂移、写入非线性和硬件实例过拟合都会在实验室原型与真实部署之间形成精度鸿沟。

本文表明，在标准器件波动下，Tiny-ViT 通过 Ensemble HAT 可达到 $86.16\pm0.19\%$ 的 fresh-instance 精度；层敏感度排序与选择性保护揭示了部署风险的长尾结构；选择性终端层 KV 缓存为 LLM 推理中的有机 CIM 应用提供了仍需锁定的候选方向。与此同时，本文也明确了边界：动态注意力的数字执行构成模拟加速上限，系统能效受外围接口和读出精度限制，ADC 分辨率存在 6-bit 阈值，硬件实例过拟合、比例噪声、保持漂移和小数据场景的数据量下限都会限制直接部署。

总体而言，有机光电 CIM 的可行性不取决于单一材料指标或单一模型精度，而取决于器件可编程性、读出精度、训练鲁棒性、数字--模拟分工和系统调度之间的协同。本文建立的仿真框架、评估协议和层级诊断方法，为后续连接真实器件测量、外围电路建模与芯片原型验证提供了方法论基础。
